\documentclass[11pt,a4paper]{article}
\usepackage[T1]{fontenc}
\usepackage[utf8]{inputenc}
\usepackage{amsmath,amssymb,xcolor}
\usepackage[margin=25mm]{geometry}
\usepackage[colorlinks=true,urlcolor=blue]{hyperref}
\definecolor{nexorablue}{HTML}{2459B5}
\setlength{\parindent}{0pt}
\setlength{\parskip}{7pt}
\setlength{\emergencystretch}{2em}
\renewcommand{\familydefault}{\sfdefault}
\begin{document}

{\small\bfseries\color{nexorablue} NEXORA CPGE}\par
{\small MPSI / NEX-CPL-04 / Corrigé}\par
{\LARGE\bfseries Un carré à partir d’une équation quartique}\par
\vspace{0.25em}\hrule\vspace{0.6em}

\section*{Racines}
Le module vaut $2$. Les arguments sont $\pi/4+k\pi/2$ pour $k=0,1,2,3$. Les points sont régulièrement espacés sur le cercle de rayon $2$.

\section*{Aire}
Le polygone est un carré. Ses diagonales sont deux diamètres perpendiculaires de longueur $4$. Son aire vaut $4\times4/2=8$.


\end{document}
